\documentclass[11pt,a4paper]{article}

\usepackage{amsmath,amssymb,amsthm}
\usepackage{mathtools}
\usepackage{geometry}
\usepackage{hyperref}
\usepackage{booktabs}
\usepackage{xcolor}
\usepackage{listings}
\usepackage{enumitem}
\usepackage{fancyhdr}
\usepackage{tocloft}
\usepackage{microtype}
\usepackage{parskip}
\usepackage{algorithm}
\usepackage{algpseudocode}

\geometry{margin=2.5cm, headheight=14pt}
\sloppy
\tolerance=3000
\emergencystretch=3em

\hypersetup{
  colorlinks=true,
  linkcolor=blue!60!black,
  urlcolor=blue!60!black,
  pdfauthor={Aegis Quant},
  pdftitle={Callable Bond Pricer -- Methodology and Mathematical Derivation}
}

\lstset{
  basicstyle=\ttfamily\small,
  keywordstyle=\color{blue!70!black}\bfseries,
  commentstyle=\color{gray},
  breaklines=true,
  frame=single,
  backgroundcolor=\color{gray!8},
  language=[Sharp]C,
  showstringspaces=false
}

\newcommand{\dd}{\mathrm{d}}
\newcommand{\E}{\mathbb{E}}
\newcommand{\Q}{\mathbb{Q}}
\newcommand{\QT}{\mathbb{Q}^T}
\newcommand{\Qt}{\mathbb{Q}^t}
\newcommand{\Qd}{\mathbb{Q}^d}
\newcommand{\R}{\mathbb{R}}
\newcommand{\N}{\mathcal{N}}
\newcommand{\F}{\mathcal{F}}
\newcommand{\Var}{\mathrm{Var}}
\newcommand{\Cov}{\mathrm{Cov}}
\newcommand{\coderef}[1]{\texttt{\color{blue!60!black}#1}}

\theoremstyle{definition}
\newtheorem{defn}{Definition}[section]
\newtheorem{prop}{Proposition}[section]
\newtheorem{remark}{Remark}[section]
\newtheorem{assumption}{Assumption}[section]

\pagestyle{fancy}
\fancyhf{}
\rhead{\textit{Aegis Quant -- Callable Bond Pricer}}
\lhead{\leftmark}
\cfoot{\thepage}

\title{\textbf{Callable Bond Pricer}\\[0.4em]
       \large Two-Factor Hull-White (G2++) Methodology,\\
       Mathematical Derivation, and Longstaff-Schwartz Implementation}
\author{Aegis Quant Library --- \texttt{MCPricer.Bonds}}
\date{\today}

\begin{document}
\maketitle
\tableofcontents
\newpage

%%==========================================================================
\section{Introduction and Scope}
\label{sec:intro}
%%==========================================================================

This document is the complete mathematical and methodological reference for
the callable bond pricer implemented in \coderef{MCPricer/Bonds}. It covers,
in full derivation:

\begin{itemize}[nosep]
  \item the two-factor Hull-White (``G2++'') short-rate model of
        Brigo \& Mercurio [\ref{brigo_mercurio}], including how it is fitted
        to an arbitrary input discount curve;
  \item the closed-form zero-coupon bond price $P(t,T)$ and the closed-form
        European option on a zero-coupon bond (ZBO), used both as the
        calibration pricing formula and as an analytical benchmark;
  \item the nonlinear least-squares calibration procedure that fits the five
        free model parameters to market ZBO prices;
  \item the Monte Carlo simulation scheme (exact factor transition,
        trapezoidal short-rate discounting);
  \item the Longstaff \& Schwartz [\ref{longstaff_schwartz}] regression-based
        valuation of the issuer's early-redemption (call) option; and
  \item the validation methodology and known numerical limitations.
\end{itemize}

Every formula below is cross-referenced to the C\# source file and method
that implements it (Section~\ref{sec:mapping} collects the full map). The
implementation lives entirely in namespace \coderef{MCPricer.Bonds}: model
(\coderef{G2ppParameters.cs}, \coderef{G2ppAnalytics.cs}), calibration
(\coderef{G2ppCalibrator.cs}, \coderef{NelderMead.cs}), and the product and
pricer (\coderef{CallableBondDefinition.cs},
\coderef{CallableBondMCPricer.cs}), with tests in
\coderef{MCPricer.Tests/Bonds}.

\subsection{Product being priced}

A \emph{callable bond} pays fixed coupons $c\cdot N/m$ at $m$ dates per year
until maturity $T_M$, at which point the principal $N$ is redeemed. In
addition, the issuer holds the right, at each date in a pre-specified
\emph{call schedule} $\{t_1^c,\dots,t_n^c\}\subset(0,T_M)$, to redeem the
bond early by paying $K_i\cdot N$ (the call price, a fraction of par) instead
of the remaining cashflows. This is a Bermudan-style option against the
bondholder: it can only reduce the bondholder's realized value relative to
the equivalent non-callable (``bullet'') bond, never increase it.

%%==========================================================================
\section{The Two-Factor Hull-White (G2++) Short-Rate Model}
\label{sec:g2pp}
%%==========================================================================

\subsection{Model dynamics and measure}

Under the risk-neutral measure $\Q$ (numeraire: the money-market account
$B(t)=\exp(\int_0^t r(s)\,\dd s)$), the short rate is decomposed into two
correlated Ornstein-Uhlenbeck (OU) factors plus a deterministic shift:
\begin{align}
  \dd x(t) &= -a\,x(t)\,\dd t + \sigma\,\dd W_1(t), \qquad x(0)=0, \label{eq:x-sde}\\
  \dd y(t) &= -b\,y(t)\,\dd t + \eta\,\dd W_2(t), \qquad y(0)=0, \label{eq:y-sde}\\
  \dd W_1(t)\,\dd W_2(t) &= \rho\,\dd t, \label{eq:corr}\\
  r(t) &= x(t) + y(t) + \varphi(t). \label{eq:short-rate}
\end{align}
Here $a,b>0$ are mean-reversion speeds, $\sigma,\eta\ge 0$ are factor
volatilities, $\rho\in(-1,1)$ is the instantaneous correlation between the
two Brownian drivers, and $\varphi(t)$ is a deterministic function chosen so
that the model reproduces an arbitrary input discount curve exactly
(derived in Section~\ref{sec:phi}). In code, the five free parameters
$(a,\sigma,b,\eta,\rho)$ are the immutable \coderef{G2ppParameters} record
(\coderef{A1, Sigma1, A2, Sigma2, Rho}).

\begin{remark}[Why two factors]
A one-factor Hull-White model produces perfectly correlated changes at every
point of the yield curve (all forward rates move together, scaled by
$B(a,\tau)$). Two imperfectly correlated factors ($\rho<1$) let short-end and
long-end rate volatility decouple, which a single factor cannot reproduce --
essential for calibrating simultaneously to short- and long-dated cap/swaption
volatilities.
\end{remark}

\subsection{Explicit solution of the OU factors}
\label{sec:ou-solution}

Equation~\eqref{eq:x-sde} is a linear SDE; multiplying by the integrating
factor $e^{at}$ and integrating gives the explicit (``exact'') solution
\begin{equation}
  x(t) = x(0)e^{-at} + \sigma\int_0^t e^{-a(t-u)}\,\dd W_1(u),
  \label{eq:x-solution}
\end{equation}
and analogously for $y(t)$ with $(b,\eta,W_2)$. Because the stochastic
integral in \eqref{eq:x-solution} is a Wiener integral of a deterministic
integrand, $x(t)$ is Gaussian conditional on $\F_0$, with
\[
  \E[x(t)] = x(0)e^{-at}, \qquad
  \Var[x(t)] = \sigma^2\int_0^t e^{-2a(t-u)}\,\dd u
             = \frac{\sigma^2}{2a}\bigl(1-e^{-2at}\bigr).
\]
By the same Itô-isometry argument applied jointly to $x$ and $y$ (using
$\dd W_1\dd W_2=\rho\,\dd t$),
\begin{equation}
  \Cov[x(t),y(t)] = \frac{\rho\sigma\eta}{a+b}\bigl(1-e^{-(a+b)t}\bigr).
  \label{eq:xy-cov}
\end{equation}

\subsection{Exact discretization for simulation}
\label{sec:exact-sim}

Equation~\eqref{eq:x-solution} gives the exact transition law over a step
$[t,t+\Delta]$: applying it with the lower integration limit shifted from $0$
to $t$,
\begin{equation}
  x(t+\Delta) = x(t)e^{-a\Delta}
              + \sigma\sqrt{\frac{1-e^{-2a\Delta}}{2a}}\;Z_1,
  \qquad Z_1\sim\N(0,1),
  \label{eq:exact-euler}
\end{equation}
and correspondingly for $y$ with $Z_2\sim\N(0,1)$, $\mathrm{Corr}(Z_1,Z_2)=\rho$.
Because \eqref{eq:exact-euler} is the \emph{exact} conditional law (not a
first-order Euler approximation), simulating $(x,y)$ this way introduces
\emph{zero} time-discretization bias regardless of step size $\Delta$ --
the only simulation error in the engine comes from the trapezoidal
integration of $r(t)$ used for discounting (Section~\ref{sec:mc-discount}),
not from the factor dynamics themselves. This is implemented in
\coderef{CallableBondMCPricer.SimulateForward} as
\lstinline{decay1 = exp(-A1*dt)}, \lstinline{vol1 = Sigma1*sqrt((1-exp(-2*A1*dt))/(2*A1))},
consuming the correlated $(Z_1,Z_2)$ pair directly from
\coderef{RandomSimulator.SimulationCube}, whose built-in Cholesky
factorization of the $2\times2$ correlation matrix $\bigl[\begin{smallmatrix}1&\rho\\\rho&1\end{smallmatrix}\bigr]$
supplies exactly the correlation required by \eqref{eq:corr}.

%%==========================================================================
\section{Fitting the Initial Term Structure}
\label{sec:phi}
%%==========================================================================

\subsection{Bond price as a conditional expectation}

By definition, the time-$t$ price of a zero-coupon bond maturing at $T$ is
\begin{equation}
  P(t,T) = \E^\Q\!\left[\exp\Bigl(-\int_t^T r(s)\,\dd s\Bigr)\,\Big|\,\F_t\right].
  \label{eq:bond-def}
\end{equation}
Since $r(s)=x(s)+y(s)+\varphi(s)$ and $\varphi$ is deterministic, and since
$x,y$ are jointly Gaussian conditional on $\F_t$ (Markov property of the OU
process), the exponent's integral $I:=\int_t^T r(s)\,\dd s$ is Gaussian
conditional on $\F_t$. Using $\E[e^{-I}]=\exp(-\E[I]+\tfrac12\Var[I])$ for a
Gaussian $I$,
\begin{equation}
  P(t,T) = \exp\!\left(-\int_t^T\varphi(s)\,\dd s
                        - \E\Bigl[\textstyle\int_t^T(x(s){+}y(s))\dd s\,\Big|\F_t\Bigr]
                        + \tfrac12\Var\Bigl[\textstyle\int_t^T(x(s){+}y(s))\dd s\,\Big|\F_t\Bigr]\right).
  \label{eq:bond-gaussian}
\end{equation}

\subsection{The conditional mean: the function $B(a,\tau)$}

Define $B(a,\tau) := \dfrac{1-e^{-a\tau}}{a}$. From
$\E[x(s)|\F_t]=x(t)e^{-a(s-t)}$,
\[
  \E\!\left[\int_t^T x(s)\,\dd s\,\Big|\,\F_t\right]
  = x(t)\int_t^T e^{-a(s-t)}\,\dd s
  = x(t)\,B(a,T-t),
\]
and analogously $\E[\int_t^T y\,\dd s|\F_t]=y(t)B(b,T-t)$.

\subsection{The conditional variance: the kernel $V(\tau)$}
\label{sec:v-derivation}

Write $\tau=T-t$. Substituting the exact solution \eqref{eq:x-solution}
(restarted at $t$) into the integral and swapping the order of integration
(Fubini):
\[
  \int_t^T x(s)\,\dd s = x(t)B(a,\tau)
    + \sigma\int_t^T\!\int_t^s e^{-a(s-u)}\dd W_1(u)\,\dd s
    = x(t)B(a,\tau) + \sigma\int_t^T B(a,T-u)\,\dd W_1(u),
\]
using $\int_u^T e^{-a(s-u)}\dd s = B(a,T-u)$. The stochastic part is a Wiener
integral, so by the Itô isometry (jointly for $x$ and $y$, using
$\dd W_1\dd W_2=\rho\,\dd t$):
{\small
\begin{align}
  \Var\!\left[\int_t^T(x(s){+}y(s))\dd s\right]
  &= \sigma^2\!\int_t^T\! B(a,T{-}u)^2\dd u
   + \eta^2\!\int_t^T\! B(b,T{-}u)^2\dd u
   + 2\rho\sigma\eta\!\int_t^T\! B(a,T{-}u)B(b,T{-}u)\dd u. \label{eq:var-integral}
\end{align}}
Substituting $v=T-u$ (so $v$ ranges over $[0,\tau]$), the three integrals in
\eqref{eq:var-integral} depend only on $\tau$ -- \emph{not} on $t$ and $T$
separately -- because the OU coefficients are time-homogeneous and
$x(0)=y(0)=0$. Evaluating each in closed form:
{\small
\[
  \int_0^\tau B(a,v)^2\,\dd v = \frac{1}{a^2}\int_0^\tau\bigl(1-e^{-av}\bigr)^2\dd v
  = \frac{1}{a^2}\left[\tau + \frac{2}{a}e^{-a\tau} - \frac{1}{2a}e^{-2a\tau} - \frac{3}{2a}\right],
\]
\[
  \int_0^\tau B(a,v)B(b,v)\,\dd v = \frac{1}{ab}\int_0^\tau\bigl(1-e^{-av}\bigr)\bigl(1-e^{-bv}\bigr)\dd v
  = \frac{1}{ab}\left[\tau + \frac{e^{-a\tau}-1}{a} + \frac{e^{-b\tau}-1}{b} - \frac{e^{-(a+b)\tau}-1}{a+b}\right].
\]}
Collecting terms, define the (time-homogeneous) \emph{variance kernel}
{\small
\begin{equation}
  \begin{aligned}
  V(\tau) := \;& \frac{\sigma^2}{a^2}\left[\tau + \frac{2}{a}e^{-a\tau} - \frac{1}{2a}e^{-2a\tau} - \frac{3}{2a}\right]
   + \frac{\eta^2}{b^2}\left[\tau + \frac{2}{b}e^{-b\tau} - \frac{1}{2b}e^{-2b\tau} - \frac{3}{2b}\right] \\
   & + \frac{2\rho\sigma\eta}{ab}\left[\tau + \frac{e^{-a\tau}-1}{a} + \frac{e^{-b\tau}-1}{b} - \frac{e^{-(a+b)\tau}-1}{a+b}\right].
  \end{aligned}
  \label{eq:v-function}
\end{equation}}
so that $\Var[\int_t^T(x{+}y)\dd s\,|\,\F_t] = V(T-t)$, independent of the
conditioning state. One checks directly that $V(0)=0$ (every bracket
vanishes at $\tau=0$), consistent with $P(t,t)=1$. Code:
\coderef{G2ppAnalytics.VFunction} (private) implements \eqref{eq:v-function}
verbatim.

\subsection{The deterministic shift $\varphi(t)$}
\label{sec:phi-derivation}

Substituting the mean and variance above into \eqref{eq:bond-gaussian} with
$t=0$, $x(0)=y(0)=0$:
\begin{equation}
  P(0,T) = \exp\!\left(-\int_0^T\varphi(s)\,\dd s + \tfrac12 V(T)\right).
  \label{eq:p0T}
\end{equation}
No-arbitrage requires $P(0,T)=P^M(0,T)$ (the input market curve) for
\emph{every} $T$. Solving \eqref{eq:p0T} for the integral and differentiating
both sides with respect to $T$ (Leibniz rule) gives $\varphi$ pointwise:
\begin{equation}
  \int_0^T\varphi(s)\,\dd s = -\ln P^M(0,T) + \tfrac12 V(T)
  \quad\Longrightarrow\quad
  \varphi(T) = f^M(0,T) + \tfrac12\,\frac{\dd V}{\dd T}(T),
  \label{eq:phi-integral-form}
\end{equation}
where $f^M(0,T)=-\partial_T\ln P^M(0,T)$ is the market instantaneous forward
rate. Differentiating \eqref{eq:v-function} term by term (each bracket's
derivative telescopes; e.g.\ $\frac{\dd}{\dd T}\bigl[\tfrac{2}{a}e^{-aT}-\tfrac{1}{2a}e^{-2aT}\bigr] = -2e^{-aT}+e^{-2aT}$, and adding the constant term's derivative $1$ gives $(1-e^{-aT})^2$):
\[
  \frac{\dd V}{\dd T}(T) = \sigma^2 B(a,T)^2 + \eta^2 B(b,T)^2 + 2\rho\sigma\eta\,B(a,T)B(b,T).
\]
Substituting into \eqref{eq:phi-integral-form} yields the closed form used by
the code:
\begin{equation}
  \boxed{\;\varphi(t) = f^M(0,t)
   + \frac{\sigma^2}{2a^2}\bigl(1-e^{-at}\bigr)^2
   + \frac{\eta^2}{2b^2}\bigl(1-e^{-bt}\bigr)^2
   + \frac{\rho\sigma\eta}{ab}\bigl(1-e^{-at}\bigr)\bigl(1-e^{-bt}\bigr)\;}
  \label{eq:phi-final}
\end{equation}
This is the \emph{unique} deterministic shift making the model reprice the
input curve exactly, for \emph{any} choice of $(a,\sigma,b,\eta,\rho)$ --
curve-fitting and volatility calibration are fully decoupled. Code:
\coderef{G2ppAnalytics.Phi}.

\begin{remark}[Instantaneous forward via finite difference]
$f^M(0,t)$ is not available in closed form from a piecewise-linear
zero-rate curve; \coderef{G2ppAnalytics.InstantaneousForward} approximates it
by a centered finite difference of \coderef{ZeroCurve.ForwardRate} over a
window $[t-h,t+h]$ with $h=10^{-4}$ years ($\approx$ 53 minutes) -- far
finer than any realistic curve pillar spacing, so the approximation error is
negligible next to curve interpolation error itself.
\end{remark}

\subsection{Closed-form bond price $P(t,T)$}
\label{sec:bond-price}

Combining Sections~\ref{sec:v-derivation}--\ref{sec:phi-derivation}: from
\eqref{eq:phi-integral-form},
\[
  \int_t^T\varphi(s)\,\dd s = \int_0^T\varphi\,\dd s - \int_0^t\varphi\,\dd s
  = \ln\frac{P^M(0,t)}{P^M(0,T)} + \tfrac12\bigl(V(T)-V(t)\bigr),
\]
(here $V(T):=V(0,T)$, $V(t):=V(0,t)$ by the time-homogeneity noted above).
Substituting into \eqref{eq:bond-gaussian} together with the mean/variance
results:
\begin{equation}
  \boxed{\;P(t,T) = \frac{P^M(0,T)}{P^M(0,t)}
    \exp\!\left(\tfrac12\bigl(V(T{-}t)-V(T)+V(t)\bigr)
                - B(a,T{-}t)\,x(t) - B(b,T{-}t)\,y(t)\right)\;}
  \label{eq:bond-price-final}
\end{equation}
Code: \coderef{G2ppAnalytics.ZeroCouponBondPrice}.

\begin{prop}[Sanity checks]
Formula \eqref{eq:bond-price-final} satisfies:
\begin{enumerate}[nosep]
  \item $P(0,T)=P^M(0,T)$ for $x(0)=y(0)=0$ (the exponent reduces to
        $\tfrac12(V(T)-V(T)+V(0))=0$ since $V(0)=0$): the model reproduces
        the input curve \emph{exactly}, independent of
        $(a,\sigma,b,\eta,\rho)$.
  \item $P(t,t)=1$ for any $x(t),y(t)$ (implemented as an explicit
        $\tau=0$ short-circuit to avoid $0/0$ in $B(a,0)$).
  \item $\partial P/\partial x(t) = -B(a,T-t)\,P(t,T) < 0$ and likewise for
        $y$: a higher realized short-rate factor strictly lowers the bond
        price, for any $a,\tau>0$.
  \item $P(t,T)$ is strictly decreasing in $T$ for fixed $t$ (positive
        rates $\Rightarrow$ longer maturities discount more).
\end{enumerate}
\end{prop}
These four properties are exercised directly by
\coderef{G2ppAnalyticsTests.cs}
(\texttt{ZeroCouponBondPrice\_\allowbreak AtT0\_\allowbreak MatchesMarketCurveExactly},
\texttt{...SameMaturityAsValuation\_\allowbreak IsOne},
\texttt{...PositiveStateShiftsDownwardBondPrice},
\texttt{...MonotonicallyDecreasingInMaturity}).

%%==========================================================================
\section{European Option on a Zero-Coupon Bond (ZBO)}
\label{sec:zbo}
%%==========================================================================

The ZBO is the model's native calibration instrument: caps, floors, and (via
standard approximations) swaptions all decompose into portfolios of options
of this type, and -- because the model is Gaussian -- its price has the same
closed form as a Black-76 option on a forward.

\subsection{Setup}

Consider a European call, expiring at $t$, on a zero-coupon bond of face
value $N$ maturing at $T>t$, struck at $X$:
\[
  \text{payoff at } t: \quad N\cdot\bigl(P(t,T)-X/N\bigr)^+ \;=\; \bigl(N\,P(t,T)-X\bigr)^+.
\]
Its time-0 price is $\E^\Q\bigl[B(t)^{-1}(N P(t,T)-X)^+\bigr]$.

\subsection{Change to the $t$-forward measure}

Switch numeraire from the money-market account $B$ to the $t$-maturity bond
$P(\cdot,t)$, defining the $t$-forward measure $\Qt$ (Section~2 of the
companion theory document, \coderef{theory.tex}, develops the
change-of-numeraire theorem in general; here $T$ there $=t$ here). Then
\[
  \text{ZBC}(t,T,X,N) = P(0,t)\;\E^{\Qt}\!\bigl[(N P(t,T)-X)^+\bigr].
\]
Girsanov's theorem states that changing measure from $\Q$ to $\Qt$ shifts the
\emph{drift} of $x(t),y(t)$ but leaves their \emph{volatility structure}
(quadratic variation) unchanged -- $x(t),y(t)$ remain jointly Gaussian under
$\Qt$ with the \emph{same} variance/covariance as under $\Q$, only a
different (in general, non-zero) mean. Since $\ln P(t,T)$ from
\eqref{eq:bond-price-final} is an \emph{affine} function of $(x(t),y(t))$,
it too remains Gaussian under $\Qt$, and its \emph{variance} under $\Qt$
equals its variance under $\Q$ conditional on $\F_0$.

\subsection{The variance $\Sigma^2(t,T)$}
\label{sec:sigma2-derivation}

From \eqref{eq:bond-price-final}, conditional on $\F_0$,
\[
  \ln P(t,T) = \text{(deterministic)} - B(a,\tau)x(t) - B(b,\tau)y(t), \qquad \tau=T-t,
\]
so, using the OU marginal variance/covariance from
Section~\ref{sec:ou-solution} (evaluated under $\Q$, and by the Girsanov
argument above equally valid under $\Qt$):
{\small
\begin{align}
  \Sigma^2(t,T) &:= \Var\bigl[\ln P(t,T)\bigr]
   = B(a,\tau)^2\Var[x(t)] + B(b,\tau)^2\Var[y(t)] + 2B(a,\tau)B(b,\tau)\Cov[x(t),y(t)] \notag\\
   &= \frac{\sigma^2}{2a^3}\bigl(1-e^{-a\tau}\bigr)^2\bigl(1-e^{-2at}\bigr)
    + \frac{\eta^2}{2b^3}\bigl(1-e^{-b\tau}\bigr)^2\bigl(1-e^{-2bt}\bigr) \notag\\
   &\qquad + \frac{2\rho\sigma\eta}{ab(a+b)}\bigl(1-e^{-a\tau}\bigr)\bigl(1-e^{-b\tau}\bigr)\bigl(1-e^{-(a+b)t}\bigr).
  \label{eq:sigma2}
\end{align}}
(Substitute $\Var[x(t)]=\tfrac{\sigma^2}{2a}(1-e^{-2at})$ and
$B(a,\tau)^2=\tfrac{1}{a^2}(1-e^{-a\tau})^2$ to see the first term collapses
to $\tfrac{\sigma^2}{2a^3}(1-e^{-a\tau})^2(1-e^{-2at})$; the cross term
follows from \eqref{eq:xy-cov} the same way.) Code:
\coderef{G2ppAnalytics.ZeroCouponBondOption} (inline, as \texttt{sigmaSq}).

\subsection{Closed-form price}

Because $\ln P(t,T)$ is Gaussian under $\Qt$ with known variance
\eqref{eq:sigma2}, and because the forward bond price
$F(0;t,T):=P(0,T)/P(0,t)$ is by construction a $\Qt$-martingale (the
defining property of the forward measure), the expectation
$\E^{\Qt}[(NP(t,T)-X)^+]$ is exactly the Black-76 formula applied to the
lognormal forward $N\cdot F(0;t,T)$ with total (non-annualized) volatility
$\Sigma(t,T)=\sqrt{\Sigma^2(t,T)}$:
\begin{equation}
  \boxed{\;
  \begin{aligned}
    \text{ZBC}(t,T,X,N) &= N\,P(0,T)\,\Phi(h_1) - X\,P(0,t)\,\Phi(h_2), \\
    \text{ZBP}(t,T,X,N) &= X\,P(0,t)\,\Phi(-h_2) - N\,P(0,T)\,\Phi(-h_1), \\
    h_1 &= \frac{1}{\Sigma}\ln\!\frac{N\,P(0,T)}{X\,P(0,t)} + \frac{\Sigma}{2}, \qquad
    h_2 = h_1 - \Sigma.
  \end{aligned}\;}
  \label{eq:zbo-formula}
\end{equation}
This is structurally identical to Jamshidian's [\ref{jamshidian1989}]
single-factor Hull-White bond-option formula, generalized to the
two-correlated-factor variance \eqref{eq:sigma2}; see Brigo \& Mercurio
[\ref{brigo_mercurio}], \S4.2. Code:
\coderef{G2ppAnalytics.ZeroCouponBondOption}, using
\coderef{AnalyticalPricers.NormalDistribution.Cdf} for $\Phi$.

\subsection{Edge cases and identities}

\begin{itemize}[nosep]
  \item \textbf{Zero-vol limit} ($\Sigma\to0$): the code returns the
        discounted intrinsic value $P(0,t)\cdot(F-X/N\cdot N)^+$ directly
        (bypassing $\log(\cdot)/\Sigma$, which is a $0/0$ singularity),
        matching the deterministic forward price
        $F=N\,P(0,T)/P(0,t)$. Verified in
        \texttt{ZboCall\_\allowbreak NearZeroVol\_\allowbreak MatchesDiscountedIntrinsic}.
  \item \textbf{Deep ITM / deep OTM}: as $X\to0$ the call price
        $\to N P(0,T)-XP(0,t)$ (pure discounted forward minus discounted
        strike); as $X\to\infty$ it $\to0$. Verified in
        \texttt{Zbo\_\allowbreak DeepItmCall\_\allowbreak ApproachesDiscountedForwardMinusStrike}
        and \texttt{Zbo\_\allowbreak DeepOtmCall\_\allowbreak IsNearZero}.
  \item \textbf{Put-call parity} (model-free, from
        $(\cdot)^+-(-\cdot)^+=\cdot$): $\text{ZBC}-\text{ZBP} = N P(0,T) - X P(0,t)$,
        verified to machine precision in \texttt{Zbo\_\allowbreak PutCallParity\_\allowbreak HoldsExactly}.
\end{itemize}

%%==========================================================================
\section{Calibration Methodology}
\label{sec:calibration}
%%==========================================================================

\subsection{Calibration instruments}

A calibration instrument (\coderef{G2ppCalibrationInstrument}) is a tuple
$(t_i,T_i,X_i,N_i,\text{isCall}_i,\text{MarketPrice}_i)$: a European
zero-coupon-bond option with a known market price. In practice these would
be derived from cap/floor or swaption quotes; that market-convention
conversion is outside the scope of this pricer, which calibrates directly
to ZBO prices.

\subsection{Objective function}

Given a candidate parameter vector $\theta=(a,\sigma,b,\eta,\rho)$, define
the model price $\hat{P}_i(\theta)$ via \eqref{eq:zbo-formula} and minimize
the mean squared \emph{relative} pricing error:
\begin{equation}
  \text{RMSE}(\theta)^2 = \frac{1}{n}\sum_{i=1}^n
    \left(\frac{\hat{P}_i(\theta)-\text{MarketPrice}_i}{\text{MarketPrice}_i}\right)^2.
  \label{eq:calib-objective}
\end{equation}
Relative errors are used (rather than absolute) so that instruments of very
different price scales -- short-dated vs.\ long-dated options -- contribute
comparably; an absolute-error objective would be dominated by the largest
premium in the instrument set. Code: \coderef{G2ppCalibrator.Calibrate},
local function \texttt{Objective}.

\subsection{Unconstrained reparameterization}

$\theta$ is constrained ($a,b>0$; $\sigma,\eta\ge0$; $\rho\in(-1,1)$), but
the optimizer (Section~\ref{sec:nelder-mead}) is unconstrained. The search is
performed over $u\in\R^5$ via the bijection
\begin{equation}
  a=e^{u_1},\quad \sigma=e^{u_2},\quad b=e^{u_3},\quad \eta=e^{u_4},\quad
  \rho=\tanh(u_5),
  \label{eq:reparam}
\end{equation}
so that every point visited by the simplex maps to a valid
\coderef{G2ppParameters} instance; no boundary-handling logic is needed
inside the objective.

\subsection{Nelder-Mead simplex minimization}
\label{sec:nelder-mead}

The optimizer is the classical derivative-free downhill simplex method of
Nelder \& Mead [\ref{nelder_mead}], chosen to avoid an external optimization
dependency (the model has only $5$ parameters and the objective is cheap
to evaluate in closed form, so gradient information is unnecessary).
Standard textbook coefficients are used throughout: reflection $\alpha=1$,
expansion $\gamma=2$, contraction $\rho_c=0.5$, shrink $\sigma_s=0.5$.

\begin{algorithm}[H]
\caption{Nelder-Mead (as implemented in \coderef{NelderMead.Minimize})}
\begin{algorithmic}[1]
\State Build an initial simplex of $n{+}1$ vertices in $\R^n$ from the guess $u_0$.
\While{iterations $<$ max \textbf{and} spread of vertex values $\ge$ tol}
  \State Sort vertices by objective value; let $\bar u$ = centroid of all but the worst.
  \State \textbf{Reflect}: $u_r = \bar u + \alpha(\bar u - u_{\text{worst}})$.
  \If{$f(u_r) < f(u_{\text{best}})$}
    \State \textbf{Expand}: $u_e = \bar u + \gamma(\bar u - u_{\text{worst}})$; keep the better of $u_r,u_e$.
  \ElsIf{$f(u_r) < f(u_{\text{second-worst}})$}
    \State Accept $u_r$.
  \Else
    \State \textbf{Contract}: $u_c = \bar u + \rho_c(u_{\text{worst}} - \bar u)$.
    \If{$f(u_c) < f(u_{\text{worst}})$} accept $u_c$
    \Else{} \textbf{shrink} every vertex toward $u_{\text{best}}$ by $\sigma_s$
    \EndIf
  \EndIf
\EndWhile
\State \Return best vertex, its value, and a \texttt{Converged} flag.
\end{algorithmic}
\end{algorithm}

\subsection{Overfitting guard and failure handling}

The model has $5$ free parameters; \coderef{G2ppCalibrator.Calibrate}
rejects (throws \coderef{ArgumentException}) any request with fewer than $5$
instruments -- an underdetermined least-squares fit is curve-fitting noise,
not calibration (see the quant-pricing review checklist item
``calibration overfitting''). After optimization, if
$\text{RMSE}(\hat\theta) > $ \texttt{maxAcceptableRmse} (default $10^{-3}$),
the method throws \coderef{InvalidOperationException} rather than returning
a poorly-fitted parameter set -- consistent with the library-wide rule of
\emph{no silent fallback on calibration failure}.

\subsection{Round-trip validation}

The correct sanity check for a nonlinear least-squares fit to a handful of
instruments is not that the raw parameter vector $\hat\theta$ matches a
``true'' generating $\theta$ bit-for-bit (the fit need not be unique), but
that \emph{the fitted model reprices the calibration instruments}. This is
exactly what \coderef{G2ppCalibratorTests.Calibrate\_\allowbreak RecoversInstrumentPrices\_\allowbreak WithinTightTolerance}
checks: synthetic ZBO prices are generated from a known $\theta_{\text{true}}$
across a 7-instrument expiry/tenor/moneyness grid, then
\coderef{G2ppCalibrator.Calibrate} is run on those prices alone, and every
instrument is re-priced under $\hat\theta$ to relative error $<10^{-3}$.

%%==========================================================================
\section{Monte Carlo Simulation of the Model}
\label{sec:mc}
%%==========================================================================

\subsection{Pricing identity}

Under the money-market numeraire, the time-0 price of any payoff stream
$\{CF_i \text{ at } t_i\}$ is
\begin{equation}
  V(0) = \E^\Q\!\left[\sum_i \frac{CF_i}{B(t_i)}\right],
  \qquad B(t)=\exp\Bigl(\int_0^t r(s)\,\dd s\Bigr).
  \label{eq:mc-price}
\end{equation}
\coderef{CallableBondMCPricer} simulates $(x,y)$ paths exactly
(Section~\ref{sec:exact-sim}), builds $r(t_i)=x(t_i)+y(t_i)+\varphi(t_i)$
on a fine uniform time grid, and accumulates $B(t_i)^{-1}$ path-by-path.

\subsection{Fine grid and $\varphi$ precomputation}

Let $T_M$ be the bond's maturity (in years from valuation) and let $n$ be the
number of steps in the supplied \coderef{SimulationCube} (its \texttt{Steps}
property); $\Delta t = T_M/n$. $\varphi(i\Delta t)$, $i=0,\dots,n$, is
precomputed once via \eqref{eq:phi-final} (shared across all paths --
\coderef{CallableBondMCPricer} constructor, field \texttt{\_\allowbreak phi}).

\subsection{Discounting via trapezoidal integration}
\label{sec:mc-discount}

The continuous integral $\int_0^{t}r(s)\,\dd s$ in $B(t)^{-1}$ is approximated
along each simulated path by the trapezoidal rule on the fine grid:
\begin{equation}
  B(t_i)^{-1} \approx \exp\!\left(-\sum_{k=1}^{i}\tfrac12\bigl(r(t_{k-1})+r(t_k)\bigr)\Delta t\right).
  \label{eq:trapezoidal}
\end{equation}
The trapezoidal rule is exact for a piecewise-linear integrand and has local
error $O(\Delta t^3)$ (global $O(\Delta t^2)$) for a smooth integrand; unlike
the factor simulation (Section~\ref{sec:exact-sim}), this step is the
\emph{only} source of time-discretization bias in the engine, and it
vanishes as $\Delta t\to0$ (Section~\ref{sec:limitations} quantifies this
in a deterministic setting where it can be measured exactly).

\subsection{Correlated increments}

The two driving Brownians in \eqref{eq:exact-euler} are supplied by
\coderef{RandomSimulator.SimulationCube.Generate(paths, steps, assets{=}2, correlation)},
where \texttt{correlation}$=\bigl[\begin{smallmatrix}1&\rho\\\rho&1\end{smallmatrix}\bigr]$.
The constructor of \coderef{CallableBondMCPricer} validates that
\texttt{cube.Correlation[0,1]} matches \coderef{G2ppParameters.Rho} to
$10^{-9}$, since the cube's own Cholesky factorization -- not any code in
the bond pricer -- is what actually correlates $Z_1,Z_2$.

%%==========================================================================
\section{The Callable Bond Product}
\label{sec:product}
%%==========================================================================

\subsection{Cashflow schedule}

\coderef{CallableBondDefinition.BuildCashflowSchedule} generates coupon
dates by stepping backward from the maturity date in
$12/m$-month increments ($m=$ \coderef{CouponsPerYear}), retaining only
dates strictly after the valuation date, using Act/365 year fractions
consistent with \coderef{ZeroCurve}. The final scheduled date carries both
the last coupon and the principal redemption $N$.

\subsection{Call schedule economics}

Each call date $t_k^c$ carries a call price $K_k$ (fraction of par). From the
\emph{bondholder's} perspective, being called at $t_k^c$ means receiving
$K_k N$ instead of the (higher, in expectation, whenever the issuer chooses
to call) remaining cashflows -- i.e.\ the bondholder is implicitly short a
Bermudan option to the issuer. Formally, if $C(t_k^c)$ denotes the
market value of continuing to hold the bond at $t_k^c$, the issuer calls
iff $K_kN < C(t_k^c)$ (calling is cheaper for the issuer than continuing to
owe the future cashflows), so the bondholder's realized value at $t_k^c$ is
\begin{equation}
  V(t_k^{c}) = \min\bigl(K_k N,\; C(t_k^c)\bigr).
  \label{eq:call-decision}
\end{equation}
This is a no-arbitrage bound: $V(t_k^c)\le C(t_k^c)$ always, hence a
callable bond can never be worth more than the corresponding bullet bond
(verified in \texttt{CallableBond\_\allowbreak NeverWorthMoreThanBulletBond}).

\subsection{Merged event grid}

\coderef{CallableBondMCPricer.BuildEventGrid} merges the coupon schedule and
the call schedule into a single chronologically sorted sequence of
\emph{event nodes}, each carrying (year-fraction, cashflow amount,
is-callable flag, call price). A date that is both a coupon date and a call
date (the typical convention -- issuers call on coupon dates) is
represented by a single node with both fields populated; this generalizes
without special-casing. Each event's year-fraction is snapshotted to the
nearest fine-grid index (Section~\ref{sec:discretization-bias} discusses the
resulting, intentionally small, discretization error).

%%==========================================================================
\section{Longstaff-Schwartz Valuation of the Embedded Call}
\label{sec:lsm}
%%==========================================================================

\subsection{Dynamic-programming formulation}

Let $t_1<\dots<t_M$ be the merged event dates ($t_M=T_M$, maturity). Define,
for each event index $k$ working \emph{backward}, the bondholder's value
process $V_k$ (a random variable indexed by simulation path) satisfying the
backward recursion
\begin{equation}
  V_M = \text{Coupon}_M + N, \qquad
  V_k = \begin{cases}
    \text{Coupon}_k + \min\bigl(K_kN,\;C_k\bigr) & \text{$t_k$ callable},\\[2pt]
    \text{Coupon}_k + C_k & \text{otherwise},
  \end{cases}
  \label{eq:bellman}
\end{equation}
where $C_k = \E^\Q\bigl[V_{k+1}\cdot B(t_k)/B(t_{k+1})\,\big|\,\F_{t_k}\bigr]$
is the continuation value discounted to $t_k$. This is a standard optimal
stopping / Bellman recursion; the price is $\E[V_1\cdot B(0)/B(t_1)]$
carried all the way back to $t=0$.

\subsection{Why regression: avoiding look-ahead bias}

$C_k$ is a conditional expectation given the Markov state $(x(t_k),y(t_k))$
-- it cannot be computed exactly from a finite set of simulated paths
without knowing the future optimal exercise policy, which is exactly what
is being determined. Longstaff \& Schwartz [\ref{longstaff_schwartz}]'s
insight is to \emph{estimate} $C_k$ by least-squares regression of the
(already-known, from the backward pass so far) realized continuation values
on a polynomial basis in the state, and to use this regression \emph{only}
to decide whether to exercise -- never as the value carried backward when
not exercising, which instead uses the actual realized (simulated)
continuation. This avoids the look-ahead bias that would result from
using an in-sample-optimal, path-omniscient exercise rule.

\subsection{Basis functions}

The regression basis used is the standard low-order polynomial basis for a
two-factor Gaussian short-rate model (matching the low-order polynomial
recommendation of [\ref{longstaff_schwartz}], \S3, adapted to two state
variables):
\begin{equation}
  \psi(x,y) = \bigl(1,\; x,\; y,\; x^2,\; y^2,\; xy\bigr) \in \R^6.
  \label{eq:basis}
\end{equation}
At each callable event $k$, ordinary least squares fits
$\hat{C}_k(x,y)=\beta\cdot\psi(x,y)$ to the cross-sectional sample
$\{(x_p(t_k), y_p(t_k), \text{pathPV}_p)\}_{p=1}^{\text{Paths}}$ by solving
the normal equations $\Psi^\top\Psi\,\beta = \Psi^\top v$ via Gaussian
elimination with partial pivoting on the resulting $6\times6$ system
(\coderef{CallableBondMCPricer.RegressContinuation},
\coderef{SolveLinearSystem}) -- small enough that no external linear-algebra
dependency is warranted.

\subsection{Reformulation in PV$_0$ units: the implemented algorithm}
\label{sec:pv0-reformulation}

Rather than carrying values forward in ``value at $t_k$'' units and
discounting stepwise (the textbook LSM presentation), the implementation
carries everything in \emph{present-value-at-0} (PV$_0$) units throughout,
which removes an entire bookkeeping layer without changing the decision:
scaling both sides of a comparison by the same positive discount factor
$D(0,t_k)$ does not change which side is larger, so
\[
  K_kN < C_k \quad\Longleftrightarrow\quad K_kN\,D(0,t_k) < C_k\,D(0,t_k),
\]
i.e.\ the call decision can equivalently compare \emph{PV$_0$-of-call-value}
against \emph{PV$_0$-of-continuation}. Let $\text{pathPV}_p$ accumulate the
PV$_0$ of all cashflows from the currently-processed event onward, under the
optimal policy determined so far (i.e.\ at later event nodes, already
processed in this backward sweep). Processing event $k$ in decreasing order:

\begin{algorithm}[H]
\caption{Backward induction at event $k$ (per path $p$), as implemented in \coderef{CallableBondMCPricer.Price}}
\begin{algorithmic}[1]
\If{node $k$ is callable}
  \State $\text{contPV}_p \gets \text{pathPV}_p$ \Comment{continuation value \emph{excluding} today's coupon}
  \State Regress $\{\text{contPV}_p\}$ on $\psi(x_p(t_k),y_p(t_k))$ $\to$ coefficients $\hat\beta$
  \For{each path $p$}
    \State $\hat{C}_p \gets \hat\beta\cdot\psi(x_p(t_k),y_p(t_k))$
    \State $\text{callValuePV}_p \gets K_k\cdot N \cdot D_p(0,t_k)$
    \If{$\text{callValuePV}_p < \hat{C}_p$}
      \State $\text{pathPV}_p \gets \text{callValuePV}_p$ \Comment{exercise: use the \emph{realized}, not estimated, call value}
    \EndIf
    \Comment{else: pathPV$_p$ unchanged $=$ its realized (simulated) continuation}
  \EndFor
\EndIf
\State $\text{pathPV}_p \gets \text{pathPV}_p + \text{Coupon}_k\cdot D_p(0,t_k)$ for every path $p$
   \Comment{coupon received unconditionally, added \emph{after} the decision}
\end{algorithmic}
\end{algorithm}

The final price and its Monte Carlo standard error are computed from
$\{\text{pathPV}_p\}_{p=1}^{\text{Paths}}$ exactly as in
\coderef{MCBasePricer.BuildResult} (sample mean, sample standard error,
$\pm1.9599639845400536\times$SE for the 95\% confidence interval).

\begin{remark}[Coupon/call coincidence]
Separating the coupon addition (unconditional, step 8) from the call
decision (steps 1--7, using the \emph{pre-coupon} continuation value)
correctly handles the common convention that a called bond still pays the
accrued coupon due on the call date: that coupon is received regardless of
the call decision, so it must not be included on either side of the
comparison in step 5 -- only the call price and the value of remaining
\emph{future} cashflows compete.
\end{remark}

%%==========================================================================
\section{Numerical Implementation Notes}
\label{sec:implementation}
%%==========================================================================

\subsection{Memory layout}

Only state at event nodes is retained across the forward pass -- not the
full fine-grid path -- via arrays \texttt{xAtEvent}, \texttt{yAtEvent},
\texttt{dAtEvent} of shape (Paths $\times$ NumEvents). NumEvents is small
(coupon frequency $\times$ years, plus call dates), so this is
$O(\text{Paths}\times\text{NumEvents})$ rather than
$O(\text{Paths}\times\text{Steps})$, letting the fine grid
(Section~\ref{sec:mc-discount}) be arbitrarily refined for discounting
accuracy without inflating memory.

\subsection{Parallelism}

The forward simulation pass (\coderef{SimulateForward}) uses
\lstinline{Parallel.For} over paths, matching the concurrency pattern used
throughout the library (\coderef{MCBasePricer}); each path writes only to
its own row of the event-snapshot matrices, so no synchronization is
needed. The backward LSM pass is inherently sequential across events (each
regression depends on the previous event's decisions) but parallel across
paths within each event.

\subsection{DV01}

\coderef{CallableBondMCPricer.ComputeDv01} bumps the discount curve by
$\pm\epsilon$ (default $1$bp) via \coderef{ZeroCurve.Shift} and reprices on
the \emph{same} \coderef{SimulationCube} (common random numbers reduce
finite-difference variance, matching the Greek methodology used by every
other pricer in the library, e.g.\ \coderef{MCBasePricer.ComputeGreeks}).
Because $\varphi(t)$ in \eqref{eq:phi-final} depends on the curve, bumping
the curve correctly perturbs \emph{both} discounting and the model drift.

%%==========================================================================
\section{Validation and Benchmarks}
\label{sec:validation}
%%==========================================================================

\begin{center}
\small
\renewcommand{\arraystretch}{1.3}
\begin{tabular}{@{}p{3.6cm}p{6.2cm}p{5.0cm}@{}}\toprule
\textbf{Check} & \textbf{What it validates} & \textbf{Test}\\\midrule
Curve reproduction & $P(0,T)=P^M(0,T)$ exactly, any $\theta$ &
  \S\ref{sec:bond-price}, \texttt{G2ppAnalyticsTests} \\
Put-call parity & Model-free ZBO identity, to machine precision &
  \texttt{Zbo\_\allowbreak PutCallParity\_\allowbreak HoldsExactly} \\
Zero-vol / deep ITM / deep OTM & ZBO collapses to discounted intrinsic &
  \texttt{ZboCall\_\allowbreak NearZeroVol\_\allowbreak ...}, \texttt{Zbo\_\allowbreak Deep*} \\
Calibration round-trip & Fitted $\hat\theta$ reprices calibration instruments &
  \texttt{Calibrate\_\allowbreak RecoversInstrumentPrices\_\allowbreak ...} \\
Bullet-bond benchmark & MC price $\approx$ closed-form curve-discount sum &
  \texttt{NoCallSchedule\_\allowbreak Matches...}, \texttt{CallPriceNeverReached\_\allowbreak Matches...} \\
No-arbitrage bound & Callable price $\le$ bullet price &
  \texttt{CallableBond\_\allowbreak NeverWorthMoreThanBulletBond} \\
Deterministic zero-vol trace & LSM decision correct with zero MC noise &
  \texttt{ZeroVol\_\allowbreak DeepItmCall\_\allowbreak ...}, \texttt{ZeroVol\_\allowbreak NoCallSchedule\_\allowbreak ...} \\
DV01 sign & Fixed-coupon bond value falls as rates rise &
  \texttt{Dv01\_\allowbreak IsNegative\_\allowbreak ForFixedCouponBond} \\
Input validation & Cube shape, correlation match, date ordering &
  \texttt{Constructor\_\allowbreak Rejects*}, \texttt{CallableBondDefinition\_\allowbreak Rejects*} \\
\bottomrule
\end{tabular}
\end{center}

All 32 tests in \coderef{MCPricer.Tests/Bonds} pass with the deterministic
seed convention used throughout the library (seed $=42$).

%%==========================================================================
\section{Known Limitations}
\label{sec:limitations}
%%==========================================================================

\begin{itemize}[nosep]
  \item \textbf{LSM in-sample bias.} Using the same simulated paths for both
        the continuation regression and the final valuation is known to
        introduce a small, generally upward, bias in American/Bermudan-style
        LSM prices -- the regression has effectively ``seen'' the payoff
        it is trying to predict. This is not corrected here; an
        (asymptotically) bias-free estimate would require an independent
        out-of-sample path set for the final valuation pass, at roughly
        double the simulation cost. Documented in the class-level comment of
        \coderef{CallableBondMCPricer}.
  \item \textbf{Event-to-grid rounding bias.} \label{sec:discretization-bias}
        Coupon dates are Act/365 calendar-day-count year fractions that in
        general do not land exactly on the uniform fine simulation grid;
        \coderef{BuildEventGrid} snapshots each event at its \emph{nearest}
        grid node, an $O(\Delta t)$ discretization error. In a controlled
        deterministic (zero-vol) test this measures at $\approx 2\times10^{-3}$
        absolute (on a par-100 bond) with a $12$-steps-per-year grid, and
        shrinks to $<10^{-3}$ with a daily grid -- consistent with an
        $O(\Delta t)$ error that vanishes as the grid is refined. This is an
        intentional engineering trade-off (uniform grid vs.\ exact per-event
        alignment) documented in \coderef{CallableBondMCPricerTests}.
  \item \textbf{Trapezoidal discounting bias.} Section~\ref{sec:mc-discount}'s
        rule is exact only for a piecewise-linear short-rate path; on a flat
        curve with $\sigma=\eta=0$ (so $r(t)\equiv\varphi(t)\equiv$ const) it
        is exact identically, since the true integral of a constant is
        trivially linear. For non-trivial $\sigma,\eta$, the bias is
        $O(\Delta t^2)$ globally and is not separately isolated from Monte
        Carlo sampling noise in the current test suite (it is dominated by
        MC standard error at the path counts used).
\end{itemize}

%%==========================================================================
\section{Code-to-Formula Mapping}
\label{sec:mapping}
%%==========================================================================

\begin{center}
\small
\renewcommand{\arraystretch}{1.3}
\begin{tabular}{@{}p{5.0cm}p{8.4cm}p{1.7cm}@{}}\toprule
\textbf{File / Member} & \textbf{Formula} & \textbf{Section}\\\midrule
\texttt{G2ppParameters.cs} & $(a,\sigma,b,\eta,\rho)$, validated & \ref{sec:g2pp} \\
\texttt{G2ppAnalytics.InstantaneousForward} & $f^M(0,t)$ (finite difference) & \ref{sec:phi-derivation} \\
\texttt{G2ppAnalytics.Phi} & $\varphi(t)$, eq.~\eqref{eq:phi-final} & \ref{sec:phi-derivation} \\
\texttt{G2ppAnalytics.VFunction} (private) & $V(\tau)$, eq.~\eqref{eq:v-function} & \ref{sec:v-derivation} \\
\texttt{G2ppAnalytics.ZeroCouponBondPrice} & $P(t,T)$, eq.~\eqref{eq:bond-price-final} & \ref{sec:bond-price} \\
\texttt{G2ppAnalytics.ZeroCouponBondOption} & ZBO, eq.~\eqref{eq:zbo-formula} & \ref{sec:zbo} \\
\texttt{G2ppCalibrationInstrument} & Calibration target tuple & \ref{sec:calibration} \\
\texttt{G2ppCalibrator.Calibrate} & RMSE objective eq.~\eqref{eq:calib-objective}, reparam.\ eq.~\eqref{eq:reparam} & \ref{sec:calibration} \\
\texttt{NelderMead.Minimize} & Downhill simplex & \ref{sec:nelder-mead} \\
\texttt{CallableBondDefinition.\allowbreak BuildCashflowSchedule} & Coupon schedule generation & \ref{sec:product} \\
\texttt{CallableBondMCPricer} ctor & Curve/cube validation, $\varphi$ precompute & \ref{sec:mc} \\
\texttt{CallableBondMCPricer.\allowbreak SimulateForward} & Exact OU step eq.~\eqref{eq:exact-euler}; trapezoidal $B(t)^{-1}$ eq.~\eqref{eq:trapezoidal} & \ref{sec:exact-sim}, \ref{sec:mc-discount} \\
\texttt{CallableBondMCPricer.\allowbreak BuildEventGrid} & Merged coupon/call event timeline & \ref{sec:product} \\
\texttt{CallableBondMCPricer.\allowbreak RegressContinuation} & OLS on basis eq.~\eqref{eq:basis} & \ref{sec:lsm} \\
\texttt{CallableBondMCPricer.\allowbreak SolveLinearSystem} & Gaussian elimination (partial pivot) & \ref{sec:lsm} \\
\texttt{CallableBondMCPricer.Price} & Backward induction eq.~\eqref{eq:bellman}, \S\ref{sec:pv0-reformulation} & \ref{sec:lsm} \\
\texttt{CallableBondMCPricer.\allowbreak AnalyticalBulletBondPrice} & Model-independent benchmark & \ref{sec:validation} \\
\texttt{CallableBondMCPricer.\allowbreak ComputeDv01} & Curve-bump Greek & \ref{sec:implementation} \\
\bottomrule
\end{tabular}
\end{center}

%%==========================================================================
\section*{References}
%%==========================================================================
\addcontentsline{toc}{section}{References}

\begin{enumerate}[label={[\arabic*]}]
  \item\label{brigo_mercurio} Brigo, D.\ and Mercurio, F.\ (2006).
        \textit{Interest Rate Models -- Theory and Practice}, 2nd ed.\
        Springer Finance. \S3 (one-factor Hull-White / Vasicek) and \S4.2
        (the two-factor additive Gaussian model, G2++): dynamics, the
        curve-fitting shift $\varphi(t)$, the zero-coupon bond price, and
        the option-on-a-zero-coupon-bond formula.
  \item\label{longstaff_schwartz} Longstaff, F.A.\ and Schwartz, E.S.\ (2001).
        ``Valuing American options by simulation: a simple least-squares
        approach.'' \textit{Review of Financial Studies} 14(1): 113--147.
  \item\label{nelder_mead} Nelder, J.A.\ and Mead, R.\ (1965).
        ``A simplex method for function minimization.''
        \textit{The Computer Journal} 7(4): 308--313.
  \item\label{jamshidian1989} Jamshidian, F.\ (1989).
        ``An exact bond option formula.''
        \textit{The Journal of Finance} 44(1): 205--209.
  \item\label{vasicek1977} Vasicek, O.\ (1977).
        ``An equilibrium characterization of the term structure.''
        \textit{Journal of Financial Economics} 5(2): 177--188.
  \item\label{hull_white1994} Hull, J.\ and White, A.\ (1994).
        ``Numerical procedures for implementing term structure models I:
        single-factor models.'' \textit{Journal of Derivatives} 2(1): 7--16.
\end{enumerate}

\end{document}
